\chapter{Lavori futuri}
\label{cap:28-lavori-futuri}

Le direzioni che seguono sono ordinate per rapporto fra valore atteso e costo,
e ciascuna \`e legata a un vincolo del \cref{cap:03-requisiti-principi} o a una
questione lasciata aperta dal \cref{cap:27-discussione}.

\section{Chiudere le questioni aperte}

Sono gli esperimenti che il progetto pu\`o eseguire con l'apparato che gi\`a
possiede, e che risolverebbero affermazioni oggi non stabilite.

\begin{description}[style=nextline, leftmargin=0pt]
  \item[\textbf{F1. Ripetere il confronto \code{SOJOURN}.}] Una sola esecuzione
  per modo non stabilisce una differenza del 10--15\%. Con cinque esecuzioni per
  modo la differenza diventerebbe difendibile o svanirebbe. \`E l'intervento a
  costo pi\`u basso e valore pi\`u alto dell'elenco.

  \item[\textbf{F2. Osservare \code{SOJOURN} fuori dal sovraccarico.}] Il modo
  \`e progettato per trovare un ginocchio, e ogni sua esecuzione \`e stata a
  capacit\`a o oltre, dove il ginocchio non esiste. Un carico a ciclo aperto con
  picco sotto la capacit\`a misurata direbbe se il modello sceglie limiti sensati
  quando ha margine.

  \item[\textbf{F3. Esercitare i pesi di \code{BUDGETED}.}] Due funzioni con SLO
  e pesi diversi, e la verifica che i limiti concessi seguano i pesi quando il
  budget \`e scarso e li ignorino quando \`e abbondante. \`E la validazione del
  contributo dichiarato che oggi manca del tutto.

  \item[\textbf{F4. Tarare la soglia di guardia al 70\%.}] \`E stata scelta, non
  misurata, e sul profilo a ciclo chiuso ha scavalcato completamente il modello.
  Una scansione su alcuni valori sotto arrivi aperti mostrerebbe se il valore
  conta.

  \item[\textbf{F5. Isolare il meccanismo del cross-talk.}] Il
  \cref{cap:25-valutazione-concorrenza} ha escluso la contesa di CPU e ha
  indicato il control plane condiviso. I candidati --- thread dello scheduler,
  event loop, strutture concorrenti delle code --- sono distinguibili per
  profilazione.
\end{description}

\section{Rimuovere il vincolo che non ha pagato}

Il \cref{cap:27-discussione} ha concluso che l'assenza di autenticazione \`e la
sola delle tre rinunce fondamentali il cui beneficio sperimentale sia
trascurabile, e la sola che impedisca l'uso fuori da un laboratorio.

Introdurre un'autenticazione a token sull'API, con verifica sul percorso prima
della risoluzione del nome, \`e un intervento contenuto. Il costo che va misurato
--- e che va misurato \emph{prima} di dichiarare l'intervento innocuo --- \`e il
suo contributo alla latenza sul percorso caldo. \`E esattamente il tipo di
domanda che questa piattaforma \`e attrezzata a rispondere.

\section{Estendere il controllo della concorrenza}

\subsection{Un ambito globale per \code{SOJOURN}}

Le due innovazioni della Parte III sono oggi disgiunte: \code{BUDGETED} decide
globalmente ma dal tempo di servizio; \code{SOJOURN} decide dal sojourn ma per
funzione. La combinazione naturale --- fabbisogno stimato dal sojourn, allocato
da un budget di piattaforma con equit\`a max-min --- non esiste, e sarebbe la
sintesi dei due risultati.

Il \cref{cap:12-concurrency-control} ha osservato che l'interfaccia dei
controllori non \`e uniforme proprio perch\'e un modo ad ambito globale non si
riconduce alla firma di uno per funzione. Quella asimmetria \`e il primo
ostacolo da rimuovere.

\subsection{Il tetto al ginocchio per \code{SOJOURN}}

Il commento nel codice registra la scelta e la sua condizione di revisione:
\emph{``no explicit knee cap --- the Little term already asks for no more than
the load needs, and the configured ceiling bounds the rest. Add one if a
function is seen climbing while its throughput stays flat.''} \`E un'osservazione
che F2 potrebbe produrre.

\subsection{Un minimo su finestra al posto delle derive}

Tre componenti mantengono estremi storici con un decadimento forfettario:
la baseline del modo adattivo (0{,}2\% per tick), e i due estremi dello
stimatore a budget (1\% per tick). Ciascuno \`e commentato come ``la versione
economica'' di un minimo su finestra con scadenza, che \`e la costruzione che
BBR e CoDel~\cite{nichols2012codel} usano.

\section{Copertura della combinatoria modulare}

Lo spazio delle configurazioni ha $2^9$ punti e la suite ne verifica due
(\cref{cap:16-moduli-diagnostici}). Una verifica di costruibilit\`a e
avviabilit\`a su tutte le combinazioni valide \`e realizzabile con gli strumenti
presenti, e coprirebbe la classe di difetti che il sistema modulare rende
possibile: una combinazione che compila e non si avvia.

Nella stessa direzione va la verifica di corrispondenza fra \code{openapi.yaml} e
le rotte effettivamente registrate (\cref{cap:05-contratti-dati}), e il cancello
end-to-end sugli scaffold generati (\cref{cap:19-cli-scaffolding}). I tre
interventi condividono la stessa forma: verificare artefatti dichiarati anzich\'e
soltanto codice eseguito.

\section{Offload: da segnale a modello}

Il modulo \code{offload} decide dal rifiuto della coda locale, senza alcun
modello del costo di trasferimento (\cref{cap:13-offload}). \`E una scelta
difendibile perch\'e non richiede calibrazione, e lascia sul tavolo tutto ci\`o
che la letteratura sull'offloading edge--cloud ha prodotto.

Due estensioni sono a portata:

\begin{itemize}[leftmargin=*]
  \item \textbf{Includere la latenza remota nella decisione.} Il gateway pu\`o
  misurare il round trip verso il target; confrontarlo con l'attesa stimata
  localmente trasformerebbe la decisione binaria in un confronto.
  \item \textbf{Offload preventivo.} Oggi si scarica solo dopo un rifiuto. Una
  frazione del traffico deviata \emph{prima} della congestione eviterebbe di
  raggiungerla.
\end{itemize}

Entrambe richiedono la campagna sperimentale che il
\cref{cap:13-offload} segnala come mancante: p95 end-to-end e tasso di rifiuto
con e senza offload, al variare del RTT verso il target.

\section{Alta disponibilit\`a: il vincolo da non rimuovere per primo}

\`E la direzione pi\`u ovvia e quella che raccomanderei per ultima.

Replicare il control plane richiede di replicare il registro, lo store delle
esecuzioni, le chiavi di idempotenza e le code, e di riconciliare le decisioni di
controllori che osservano stati parzialmente sovrapposti. Il
\cref{cap:03-requisiti-principi} ha argomentato che le oscillazioni cos\`i
introdotte sono un fenomeno \emph{diverso} da quelli studiati qui, e
indistinguibili da essi.

Il valore di una versione ad alta disponibilit\`a \`e in un'altra direzione: il
coordinamento fra controllori replicati \`e un problema interessante di per s\'e,
e questa piattaforma sarebbe un buon banco per studiarlo. Ma sarebbe una
piattaforma diversa, e il \cref{cap:29-conclusioni} sostiene che la sua
utilit\`a attuale \`e proporzionale a ci\`o che non fa.

\section{Sintesi}

\begin{table}[htbp]
\centering
\caption{Direzioni future, per rapporto valore/costo.}
\label{tab:futuro}
\small
\begin{tabularx}{\textwidth}{@{}lXll@{}}
\toprule
\textbf{Id} & \textbf{Intervento} & \textbf{Costo} & \textbf{Valore} \\
\midrule
F1 & Ripetere il confronto \code{SOJOURN} & basso & alto \\
F3 & Esercitare i pesi di \code{BUDGETED} & basso & alto \\
F2 & \code{SOJOURN} fuori dal sovraccarico & basso & medio \\
F5 & Isolare il meccanismo del cross-talk & medio & alto \\
--- & Autenticazione a token & basso & alto (uso) \\
--- & \code{SOJOURN} ad ambito globale & alto & alto \\
--- & Combinatoria modulare verificata & medio & medio \\
--- & Offload con modello di costo & medio & medio \\
--- & Alta disponibilit\`a & molto alto & basso per gli obiettivi attuali \\
\bottomrule
\end{tabularx}
\end{table}

Le prime quattro righe hanno una propriet\`a comune: non richiedono codice nuovo.
Sono esperimenti da eseguire su ci\`o che esiste gi\`a, e la loro presenza in cima
alla lista \`e coerente con la tesi del documento --- che il valore di questa
piattaforma stia in ci\`o che consente di misurare.
